% Mo's Algorithm on Trees

Offline queries on trees using Euler Tour (flattening the tree into an array). Set $\text{BLOCK\_SIZE} \approx \frac{2N}{\sqrt{Q}}$ for an optimal time complexity of $O((N + Q)\sqrt{N})$.

\begin{lstlisting}
const int BLOCK = 450;
const int MAXN = 100005;

struct Query {
    int l, r, lc, id;
    bool operator<(const Query& o) const {
        int b1 = l / BLOCK, b2 = o.l / BLOCK;
        if (b1 != b2) return b1 < b2;
        return (b1 & 1) ? (r < o.r) : (r > o.r);
    }
};

int n, q, timer;
int st[MAXN], en[MAXN], id[2 * MAXN], depth[MAXN], up[MAXN][20];
vector<int> adj[MAXN];
bool vis[MAXN];

void dfs(int u, int p) {
    st[u] = ++timer; id[timer] = u;
    up[u][0] = p;
    for (int i = 1; i < 20; ++i) up[u][i] = up[up[u][i-1]][i-1];
    for (int v : adj[u]) {
        if (v != p) { depth[v] = depth[u] + 1; dfs(v, u); }
    }
    en[u] = ++timer; id[timer] = u;
}

int get_lca(int u, int v) {
    if (depth[u] < depth[v]) swap(u, v);
    for (int i = 19; i >= 0; --i) {
        if (depth[u] - (1 << i) >= depth[v]) u = up[u][i];
    }
    if (u == v) return u;
    for (int i = 19; i >= 0; --i) {
        if (up[u][i] != up[v][i]) { u = up[u][i]; v = up[v][i]; }
    }
    return up[u][0];
}

void check(int u) {
    if (vis[u]) { /* remove_from_data_structure(u); */ }
    else { /* add_to_data_structure(u); */ }
    vis[u] ^= 1;
}

vector<long long> solve_mo_tree(vector<Query>& queries) {
    timer = 0; dfs(1, 1);
    for (auto& q : queries) {
        int u = q.l, v = q.r; // Here q.l and q.r store original tree node indices
        int lca = get_lca(u, v);
        if (st[u] > st[v]) swap(u, v);
        if (lca == u) {
            q.l = st[u]; q.r = st[v]; q.lc = 0;
        } else {
            q.l = en[u]; q.r = st[v]; q.lc = lca;
        }
    }
    sort(queries.begin(), queries.end());
    
    vector<long long> answers(queries.size());
    int cur_l = 1, cur_r = 0;
    for (const auto& q : queries) {
        while (cur_l > q.l) check(id[--cur_l]);
        while (cur_r < q.r) check(id[++cur_r]);
        while (cur_l < q.l) check(id[cur_l++]);
        while (cur_r > q.r) check(id[cur_r--]);
        if (q.lc) check(q.lc);
        answers[q.id] = 0; // fetch current answer here
        if (q.lc) check(q.lc);
    }
    return answers;
}
\end{lstlisting}

\begin{itemize}
    \item Convert path queries into range queries on the \lstinline|id| array size $2N$ ($1$-indexed).
    \item If LCA is one of the endpoints, the range is \lstinline|[st[u], st[v]]|. Otherwise, the range is \lstinline|[en[u], st[v]]| plus the isolated \lstinline|lca| node.
    \item The function \lstinline|check(u)| toggles the presence of node $u$, handling insertions and deletions symmetrically.
\end{itemize}
